% !TeX program = pdflatex
% =============================================================================
% Lektionsplan (Diskussionsgrundlage, NICHT das Arbeitsblatt selbst)
% Gewitter & Beginn Stromkreise -- UZH Praktikum I, Sara Romer Urban, KS im Lee, HS2026
% Klasse: 1f
% =============================================================================
\documentclass[11pt,a4paper]{article}

\usepackage[T1]{fontenc}
\usepackage[utf8]{inputenc}
\usepackage[ngerman]{babel}
\usepackage{lmodern}
\usepackage[a4paper,margin=2.2cm]{geometry}
\usepackage{array,booktabs,tabularx,xltabular}
\usepackage{enumitem}
\usepackage{xcolor}
\usepackage{amsmath}
\usepackage{siunitx}
\sisetup{per-mode=fraction,fraction-command=\frac}
\usepackage{url}
\Urlmuskip=0mu plus 1mu
\def\UrlBreaks{\do\/\do-\do_}

\definecolor{titleblue}{HTML}{183A6B}
\definecolor{accentblue}{HTML}{2E6F9E}
\definecolor{groupteal}{HTML}{1B7F72}
\definecolor{darkgray}{HTML}{4A4A4A}
\definecolor{warnred}{HTML}{9E2E2E}

\setlength{\parindent}{0pt}
\setlength{\parskip}{0.5em}

\newcommand{\Methode}[1]{{\small\color{groupteal}\textbf{#1}}}

\usepackage{titlesec}
\titleformat{\section}{\Large\bfseries\color{titleblue}}{\thesection}{0.6em}{}
\titlespacing*{\section}{0pt}{1.2em}{0.5em}

\begin{document}
\sloppy

\begin{center}
{\Huge\bfseries\color{titleblue}Lektionsplan: Gewitter \& Beginn Stromkreise}\\[0.4em]
{\large Doppellektion (2$\times$45 Min.) -- Diskussionsgrundlage}
\end{center}
\vspace{0.5em}

\noindent\textbf{Klasse:} 1f (09.09.2026)
\noindent\textbf{Lehrperson:} Loris Keller

\vspace{0.3em}
\noindent\textcolor{darkgray}{Dieses Dokument ist \textbf{keine} Reinschrift des Arbeitsblatts, sondern die
Planungsgrundlage dafür: Lernziele, Aufbau der Lektion und die gewählte
Didaktik-Methode pro Phase (aus \texttt{Methodensammlung\_Physikunterricht.pdf}).
Nach Absprache wird daraus das \texttt{.tex}-Arbeitsblatt/Aufgabenblatt lokal in
VS Code erstellt. Der Aufbau folgt dem von Loris vorgeschlagenen Ablauf: erste
Lektion Gewitter (Ladungsentstehung $\to$ elektrisches Feld, gemeinsam im
Plenum besprochen $\to$ Faradayscher Käfig mit Elektroskop-Experiment $\to$
Gruppenarbeit: SuS erstellen selbst eine Skizze und einen Erklärtext zu
Ladungsverteilung, Influenz und elektrischem Feld $\to$ Plenumsbesprechung
$\to$ Schrittspannung), zweite Lektion Beginn
Stromkreise (kurze Einführung Wasserkreislauf-Analogie und Batterie als
Ladungspumpe $\to$ \textbf{Gruppenpuzzle (D3)} zu Strom, Spannung und
Widerstand, je qualitativ mit Wasseranalogie und Alltagsbeispiel $\to$
gemeinsame Vertiefungsaufgabe). \textbf{Arbeitshypothese:} beide Hälften
bilden zusammen die eine Doppellektion \texttt{e-feld-vertiefung} vom
09.09.2026; das quantitative Ohm'sche Gesetz und die exakte Definition der
Spannung ordne ich der folgenden, separaten Doppellektion \texttt{u-i-r}
(16.09.2026) zu -- bitte kurz bestätigen oder korrigieren.}

% =========================================================================
\section*{1. Abgrenzung und Lernziele}

\textbf{Im Fokus:} zwei Anwendungen des bereits bekannten elektrischen
Feldes -- erstens Entstehung und Gefahr von Blitzen sowie der Schutz durch
den Faradayschen Käfig (PhysikLibre 12.6.11, 12.4.5/12.4.6), zweitens der
Einstieg in den Stromkreis über die Wasserkreislauf-Analogie (PhysikLibre
12.8/12.8.1/12.8.2), gefolgt von einem Gruppenpuzzle, in dem die SuS Strom
(PhysikLibre 12.6.1), Spannung und Widerstand je qualitativ mit
Wasseranalogie und Alltagsbeispiel erarbeiten und sich gegenseitig
erklären.

\textbf{Bereits bekannt (nur Repetition, keine Neueinführung):} Ladung
(gleichnamig/ungleichnamig), Coulombsche Fernwirkungskraft, elektrisches
Feld und Feldlinienregeln, Influenz, homogenes vs. inhomogenes Feld,
Feldbeispiele Punktladung/Dipol/Plattenkondensator (alle: letzte
Doppellektion, \texttt{ladung-e-feld/}).

\textbf{Bewusst nicht Teil dieser Einheit:} das Ohm'sche Gesetz und die
quantitative Beziehung $R=U/I$ sowie die exakte physikalische Definition
der Spannung als Verschiebearbeit pro Ladung ($U=W/Q$, PhysikLibre 12.5)
-- eigene Lektion \texttt{u-i-r/} (1f, 16.09.2026, eine Woche später);
Reihen-/Parallelschaltung -- eigene Lektion \texttt{schaltungen/} (1f,
23.09.2026); die exakte Berechnung der Schrittspannung (keine Formel,
Spannung/Widerstand hier noch nicht formal eingeführt -- rein
konzeptionell als Sicherheitsaspekt). Spannung selbst wird in dieser
Doppellektion bereits \textbf{qualitativ} eingeführt (Potenzialgefälle/
treibende Kraft, analog zum Höhenunterschied im Wasserkreislauf) -- nur
Formel und Rechnung folgen erst in \texttt{u-i-r/}. Auf Einheiten
(Ampere, Volt) wird bewusst ganz verzichtet, ebenso auf Aufgaben, die
mehr Vorwissen bräuchten, als die Klasse bisher hat.

\begin{itemize}[leftmargin=1.3cm,itemsep=0.3em]
  \item[LZ1] die Ladungstrennung in einer Gewitterwolke durch
  Reibungselektrizität zwischen Eiskristallen und aufsteigender warmer Luft
  erklären und beschreiben, wie daraus ein elektrisches Feld entsteht, das
  bei Überschreiten einer kritischen Feldstärke zum lawinenartigen
  Ladungsausgleich (Blitz) führt (PhysikLibre 12.6.11).
  \item[LZ2] anhand eines Experiments (Elektroskop in einem Faradayschen
  Käfig, geladener Stab von aussen angenähert) beobachten, dass sich im
  Käfig-Inneren nichts ändert, und dies erklären: ein Leiter im
  elektrostatischen Gleichgewicht ist innen feldfrei, weil Influenzladungen
  das äussere Feld exakt kompensieren (PhysikLibre 12.4.5/12.4.6).
  \item[LZ3] diese Käfig-Wirkung auf Auto und Haus übertragen und
  begründen, warum ein Auto auch trotz isolierender Gummireifen als
  Faradayscher Käfig schützt (der Blitz überspringt die Lücke zum Boden).
  \item[LZ4] Schrittspannung als Gefahr bei einem Blitzeinschlag in der
  Nähe qualitativ beschreiben (Potenzialgefälle entlang des Bodens um die
  Einschlagstelle) und daraus Sicherheitsregeln bei Gewitter ableiten, ohne
  Berechnung.
  \item[LZ5] die Grundelemente eines Stromkreises (Batterie, Schalter,
  Verbraucher) über die Wasserkreislauf-Analogie zuordnen (Pumpe
  $\hat=$ Batterie, Ventil $\hat=$ Schalter, Wasserrad $\hat=$ Verbraucher,
  z.\,B. eine LED, Wasser $\hat=$ frei bewegliche Ladungsträger, Rohre
  $\hat=$ Kabel) und ein Bild einer Batterie interpretieren, die als
  Ladungspumpe die Pole trennt und ein Potenzialgefälle erzeugt.
  \item[LZ6] (Expertenthema Strom) elektrischen Strom als die pro
  Zeiteinheit durch eine Fläche fliessende Ladung erklären
  ($I=\Delta Q/\Delta t$, ohne Einheit/Zahlenwerte), mit
  Wasseranalogie (Durchflussmenge) und Alltagsbeispiel, und dies der
  eigenen Puzzle-Gruppe vortragen.
  \item[LZ7] (Expertenthema Spannung) Spannung qualitativ als
  Potenzialgefälle/treibende Kraft der Batterie erklären, mit
  Wasseranalogie (Höhenunterschied/Druck der Pumpe) und Alltagsbeispiel
  (volle vs. fast leere Batterie, ohne Volt-Zahlen), und dies der
  eigenen Puzzle-Gruppe vortragen.
  \item[LZ8] (Expertenthema Widerstand) den elektrischen Widerstand
  qualitativ als das erklären, was den Ladungsfluss erschwert, mit
  Wasseranalogie (enges Rohr $\hat=$ hoher Widerstand) und
  Alltagsbeispiel, unter Bezug auf Leitung/Leiter, und dies der eigenen
  Puzzle-Gruppe vortragen.
  \item[LZ9] (Vertiefungsaufgabe) Strom, Spannung und Widerstand in
  einer gemeinsamen Aufgabe richtig anwenden und voneinander abgrenzen.
\end{itemize}

% =========================================================================
\section*{2. Aufbau der Doppellektion}

\subsection*{Erste Lektion: Gewitter (ca. 44 Min.)}

\begin{xltabular}{\textwidth}{@{}p{2.6cm}X p{1.2cm}@{}}
\toprule
\textbf{Phase} & \textbf{Inhalt \& Methode} & \textbf{Zeit} \\
\midrule
\endhead

Repetition &
\textbf{Ladung und elektrisches Feld.} Kurzer Rückblick auf die letzte
Doppellektion: Feldlinienregeln, Influenz, homogenes/inhomogenes Feld.
Plenumsgespräch. &
4' \\[0.4em]

Aktivierung &
\textbf{Gewitterwolke -- Einstiegsbild.} Bild einer Gewitterwolke
(\texttt{Gewitter\_Grafik\_Jung.jpg}). Frage: wie entstehen dort
überhaupt Ladungen? Plenumsgespräch, Sammeln von Vermutungen. &
3' \\[0.4em]

Theorie &
\textbf{Ladungstrennung und elektrisches Feld im Gewitter.}
Reibungselektrizität zwischen Eiskristallen und aufsteigender warmer
Luft, Oberseite der Wolke positiv, Unterseite negativ; Erde bildet bei
Gewitter den positiven Pol (\texttt{atmospheric-electric-field-static.svg}).
Bei kritischer Feldstärke: lawinenartiger Ladungsausgleich, der Blitz
(LZ1). Gemeinsam im Plenum erarbeitet, von der Lehrperson moderiert. &
6' \\[0.4em]

Übergang &
\textbf{Wie schützt man sich?} Überleitung: wenn Blitze durch ein
elektrisches Feld entstehen -- wie kann man sich davor schützen? Plenum. &
1' \\[0.4em]

Experiment &
\textbf{Elektroskop im Faradayschen Käfig.} Demonstration: ein Elektroskop
befindet sich in einem Käfig (\texttt{Elektroskop.svg},
\texttt{Faraday\_1.png}--\texttt{Faraday\_4.png}); ein geladener Stab wird
von aussen angenähert -- das Elektroskop schlägt nicht aus. Reine
Beobachtung ohne Vorhersage, da die Klasse noch keine Vorkenntnis hat, um
das Ergebnis vorherzusagen (bewusst \textbf{kein} Predict-Observe-Explain/
M4, das genau diese Vorkenntnis voraussetzen würde). Die Erklärung folgt
nicht durch die Lehrperson, sondern in der nächsten Phase durch die SuS
selbst. &
8' \\[0.4em]

Gruppenarbeit &
\textbf{Skizze und Erklärtext.} Jede Gruppe erstellt eine eigene Skizze
mit Ladungsverteilung, Influenz und elektrischem Feld sowie einen kurzen
Text, der mit diesen drei Begriffen erklärt, weshalb das Elektroskop im
Käfig nicht ausschlägt (LZ2). \Methode{M3 Whiteboarding} (angelehnt --
hier auf Papier/Poster statt auf einem Whiteboard, sonst gleiche
Struktur: Gruppenphase mit Skizze/Erklärung, danach Klassendiskussion).
Gruppenarbeit. &
11' \\[0.4em]

Plenum &
\textbf{Präsentation und Sicherung.} Gruppen zeigen und vergleichen ihre
Skizzen/Texte; Lehrperson moderiert, fragt bei rein additiver
Begriffsnennung gezielt nach dem Kausalzusammenhang nach (Influenz
$\to$ Gegenfeld $\to$ Kompensation) und sichert die korrekte Erklärung.
Anschliessend Übertragung auf Auto (\texttt{faraday-cage-car-static.jpg},
\texttt{auto\_innenraum\_blitzschlag.png}) und Haus -- die Gummireifen
spielen dabei keine Rolle, der Blitz überspringt die Lücke zum Boden
(LZ2, LZ3). Klassendiskussion, Lehrperson moderiert. &
7' \\[0.4em]

Diskussion &
\textbf{Schrittspannung.} Warum ist es bei einem nahen Blitzeinschlag
gefährlich, mit gespreizten Beinen zu stehen oder zu rennen?
(\texttt{schrittspannung.jpg}) \Methode{C3 Think-Pair-Share}, danach
Ableitung der Sicherheitsregel (Beine zusammen, in die Hocke) im Plenum
(LZ4). &
4' \\

\bottomrule
\end{xltabular}

\textcolor{darkgray}{\small Summe erste Lektion: 44 Min.}

\vspace{0.6em}
\subsection*{Zweite Lektion: Beginn Stromkreise (ca. 43 Min.)}

\begin{xltabular}{\textwidth}{@{}p{2.6cm}X p{1.2cm}@{}}
\toprule
\textbf{Phase} & \textbf{Inhalt \& Methode} & \textbf{Zeit} \\
\midrule
\endhead

Übergang &
\textbf{Vom ungeordneten zum kontrollierten Ladungsfluss.} Kurze
Überleitung: beim Blitz fliessen Ladungen unkontrolliert durch die Luft --
in einem Stromkreis fliessen sie kontrolliert durch Kabel. Plenum. &
2' \\[0.4em]

Theorie &
\textbf{Kurze Einführung: Wasserkreislauf-Analogie und Batterie.}
Grundelemente: Pumpe $\hat=$ Batterie, Ventil $\hat=$ Schalter, Wasserrad
$\hat=$ Verbraucher (z.\,B. eine LED), Wasser $\hat=$ frei bewegliche
Ladungsträger, Rohre $\hat=$ Kabel (\texttt{water-circuit2-static.png},
\texttt{water-current-circuits-static.jpg}, \texttt{batterie.png}),
daneben ein einfaches Stromkreis-Schaltbild
(\texttt{simple-circuit2-static.svg}): die Batterie trennt Ladungen und
erzeugt so ein Potenzialgefälle, das den Ladungsfluss im Kreis antreibt
(LZ5). Bewusst knapp gehalten, da Strom/Spannung/Widerstand gleich danach
vertieft im Gruppenpuzzle erarbeitet werden. Lehrervortrag. &
8' \\[0.4em]

Übung &
\textbf{Zuordnung Wasserkreislauf $\leftrightarrow$ Stromkreis.} SuS
ordnen die Elemente aus dem Wasserkreislauf-Bild den entsprechenden
Stromkreis-Bauteilen zu (LZ5). \Methode{C1 Sortieraufgabe},
Plenumsabgleich. &
6' \\[0.4em]

Expertenphase &
\textbf{Gruppenpuzzle: Strom, Spannung, Widerstand.} Stammgruppen zu
dritt bilden sich, jedes Mitglied wird Expertin/Experte für eines der
drei Themen (Karten je mit Definition, Wasseranalogie und
Alltagsbeispiel, siehe Abschnitt 3) und bereitet sich vor, es den
anderen zu erklären (LZ6, LZ7, LZ8). \Methode{D3 Partner-/
Gruppenpuzzle} -- Expertenphase, alle SuS mit demselben Thema arbeiten
zunächst zusammen. Gruppenarbeit. &
10' \\[0.4em]

Austausch &
\textbf{Gruppenpuzzle: gegenseitiges Erklären.} Neue Dreiergruppen mit
je einer Expertin/einem Experten pro Thema; reihum erklärt jede/r ihr
Thema den anderen beiden (LZ6, LZ7, LZ8). \Methode{D3 Partner-/
Gruppenpuzzle} -- Austauschphase. Gruppenarbeit. &
9' \\[0.4em]

Vertiefung &
\textbf{Gemeinsame Aufgabe.} Die Puzzle-Gruppen lösen zusammen eine
kurze, bewusst einfache Aufgabe, die Strom, Spannung und Widerstand
gemeinsam braucht (LZ9) -- empfohlen: \glqq Fehlersuche im
Stromkreis\grqq{} (eine LED leuchtet nicht -- Batterie leer? Kabel
unterbrochen? Schalter offen? -- mit der Wasseranalogie erklärt), als
Zeitreserve ein kurzes Alltagsrätsel (Taschenlampe leuchtet schwächer
bei fast leerer Batterie). Zwei konkrete Varianten in Abschnitt 3, ohne
Einheiten/Zahlenwerte und ohne Vorwissen, das über diese Lektion
hinausgeht. Ausblick: die genaue Beziehung zwischen Strom, Spannung und
Widerstand (Ohm'sches Gesetz) folgt in der nächsten Doppellektion
(\texttt{u-i-r}, 16.09.2026). Gruppenarbeit, kurzer Plenumsabgleich. &
6' \\[0.4em]

Abschluss &
\textbf{Zusammenfassung und Ausblick.} Kurze Zusammenfassung, Hinweis
auf das separate Aufgabenblatt für weitere Übung. Plenum. &
2' \\

\bottomrule
\end{xltabular}

\textcolor{darkgray}{\small Summe: 44 + 43 = 87 Min. -- 3 Min. Puffer
innerhalb der 90 Min. der Doppellektion.}

% =========================================================================
\section*{3. Expertenkarten Gruppenpuzzle (Kurzinhalt)}

Für die Expertenphase braucht jedes der drei Themen eine kurze,
möglichst einfach gehaltene Theoriekarte -- Definition, Wasseranalogie
und ein Alltagsbeispiel, dazu Bezug zu bereits Bekanntem (Ladung,
Influenz, elektrisches Feld, Leitung). Das eigentliche Kartendesign
entsteht später im Arbeitsblatt; hier die inhaltliche Kurzfassung als
Diskussionsgrundlage:

\begin{itemize}[leftmargin=1.2cm,itemsep=0.5em]
  \item \textbf{Strom (LZ6):} Definition $I=\Delta Q/\Delta t$ (keine
  Einheit, keine Zahlenwerte) (\texttt{electric-charge-and-electric-current-static.svg}).
  Wasseranalogie: die Wassermenge, die pro Zeit durch das Rohr fliesst.
  Alltagsbeispiel: ein dickeres Kabel \glqq verträgt\grqq{} mehr Strom,
  so wie ein dickeres Rohr mehr Wasser durchlässt. Bezug zu Bekanntem:
  die frei beweglichen Ladungsträger (Ladung), die durch den Leiter
  fliessen.
  \item \textbf{Spannung (LZ7):} qualitativ als Potenzialgefälle/
  treibende Kraft, die die Batterie zwischen ihren Polen erzeugt
  (\texttt{batterie.png}). Wasseranalogie: der Höhenunterschied bzw.
  Druck, den die Pumpe erzeugt -- je grösser, desto stärker \glqq
  drängt\grqq{} das Wasser durch die Rohre. Alltagsbeispiel: eine volle
  Batterie treibt stärker an als eine fast leere (bewusst ohne
  Volt-Zahlen). Bezug zu Bekanntem: das elektrische Feld, das die
  Batterie im Stromkreis erzeugt und das die Ladungsträger antreibt --
  anders als beim Faradayschen Käfig (erste Lektion) gleicht sich
  dieses Feld nie aus, weil die Batterie es ständig erneuert, solange
  der Kreis geschlossen ist.
  \item \textbf{Widerstand (LZ8):} qualitativ als das, was den
  Ladungsfluss erschwert (\texttt{resistor-static.jpg}). Wasseranalogie:
  ein enges Rohr entspricht einem hohen Widerstand. Alltagsbeispiel: ein
  Heizdraht leitet schlechter als ein dickes Kupferkabel und wird
  deshalb warm. Bezug zu Bekanntem: Leitung/Leiter -- wie leicht das
  Material die Ladungsträger hindurchlässt.
  \texttt{drude-model-static.svg} als optionale Vertiefung zur
  mikroskopischen Vorstellung, sonst für \texttt{u-i-r/} reservieren.
\end{itemize}

\textbf{Vertiefungsaufgabe (LZ9, nach dem Puzzle) -- zwei bewusst
einfach gehaltene Varianten zur Auswahl, ohne Einheiten/Zahlenwerte und
ohne Vorwissen, das über diese Lektion hinausgeht (die
\glqq Was passiert, wenn\ldots{}\grqq{}-Variante und die
Serie-/Parallelschaltung wurden deshalb gestrichen):}

\begin{itemize}[leftmargin=1.2cm,itemsep=0.4em]
  \item \textbf{(a) Fehlersuche im Stromkreis (empfohlen):} eine LED
  leuchtet nicht. Drei mögliche Ursachen -- Batterie leer, Kabel
  unterbrochen, Schalter offen -- sollen mit der Wasseranalogie erklärt
  werden (Pumpe liefert keinen Druck / Rohr unterbrochen / Ventil zu).
  Alltagsanker: ein Handy-Ladekabel, das nicht mehr lädt.
  \item \textbf{(b) Alltagsrätsel (Zeitreserve):} warum leuchtet eine
  Taschenlampe schwächer, wenn die Batterien fast leer sind? Passt
  direkt zum Alltagsbeispiel der Spannungskarte (volle vs. fast leere
  Batterie).
\end{itemize}

% =========================================================================
\section*{4. Offene Punkte zur Besprechung}

\begin{itemize}[leftmargin=1.2cm,itemsep=0.4em]
  \item \textbf{Aufteilung auf zwei Lektionen:} Ich gehe davon aus, dass
  \glqq erste Lektion\grqq{} und \glqq zweite Lektion\grqq{} die beiden
  45-Minuten-Hälften dieser einen Doppellektion \texttt{e-feld-vertiefung}
  (09.09.2026) sind, und dass das quantitative Ohm'sche Gesetz erst in
  \texttt{u-i-r} (16.09.2026) folgt. Bitte bestätigen oder korrigieren,
  falls die zweite Lektion (Stromkreise) inhaltlich bereits zu
  \texttt{u-i-r} gehören soll.
  \item \textbf{Spannung bereits qualitativ hier?} Bisher war Spannung
  komplett auf \texttt{u-i-r} verschoben; für das Gruppenpuzzle führe ich
  sie nun schon hier -- rein qualitativ, ohne Formel -- ein, da sie sich
  sonst nicht sinnvoll von Strom und Widerstand abgrenzen liesse. Bitte
  kurz bestätigen, dass das keine unerwünschte Vorwegnahme von
  \texttt{u-i-r} ist.
  \item \textbf{Schrittspannung ohne Formel:} Da Spannung und Widerstand
  an dieser Stelle noch nicht formal eingeführt sind, wird Schrittspannung
  hier rein qualitativ (Sicherheitsaspekt) behandelt. Falls stattdessen
  doch eine einfache quantitative Herangehensweise gewünscht ist, bitte
  Rücksprache -- das würde eine Vorwegnahme von Inhalten aus
  \texttt{u-i-r} bedeuten.
  \item \textbf{Elektroskop-Experiment:} Ich gehe von einem
  Demonstrationsexperiment der Lehrperson aus (nicht von SuS selbst
  durchgeführt), gefolgt von der Gruppenarbeit mit Skizze und Erklärtext.
  Bitte kurz bestätigen, ob das dem geplanten Aufbau entspricht.
  \item \textbf{Skizze/Text -- Papier oder Poster?} Ich habe die
  Gruppenarbeit an \Methode{M3 Whiteboarding} angelehnt, aber mit
  Papier/Poster statt Whiteboard, da mir kein Hinweis auf Whiteboards im
  Schulzimmer vorliegt. Falls tatsächlich Whiteboards zur Verfügung
  stehen, kann die Methode unverändert als M3 durchgeführt werden
  (Boards hochhalten statt Poster vorzeigen).
  \item \textbf{Gruppengrösse Gruppenpuzzle:} Das Puzzle setzt
  Dreiergruppen voraus (ein Thema pro Person). Falls die Klassengrösse
  nicht durch drei teilbar ist, braucht es einzelne Vierergruppen (ein
  Thema doppelt) oder eine kleine Gruppe mit nur zwei Themen -- bitte
  kurz Rücksprache, wie gross die Klasse 1f aktuell ist.
  \item \textbf{Zeitbudget:} Sowohl die Gruppenarbeit in der ersten
  Lektion (11') als auch Expertenphase (10') und Austauschphase (9') in
  der zweiten sind eher knapp bemessen. Falls das erfahrungsgemäss mehr
  Zeit braucht, ist die Übung \glqq Zuordnung Wasserkreislauf
  $\leftrightarrow$ Stromkreis\grqq{} (6', zweite Lektion) der Kandidat
  zum Kürzen oder Streichen, da dieselbe Zuordnung im Aufgabenblatt
  vertieft werden kann.
\end{itemize}

% =========================================================================
\section*{5. Anschlussdokument: Aufgabenblatt}

Wie bei den vorherigen Themen entsteht zusätzlich ein separates
Aufgabenblatt mit weiteren Übungsaufgaben zu Gewitter/Faradayschem Käfig
und zu den Grundelementen des Stromkreises, inklusive weiterer
Zuordnungs- und Verständnisaufgaben zur Wasserkreislauf-Analogie.

\end{document}
